\subsubsection{Type inference}
Type inference in general fails, if two (or more) branches fail to resolve to unifiable types.

We start a \bi{type judgement} with judgement $\vdash t :: \tau_0$, then build a derivation tree bottom-up.
Finally, apply constraints / unification to get possible types.


\paragraph{Self application}
This means that you apply a function to itself.
In Haskell, this is not typeable because there would need to be an infinite function type, but all \bi{Haskell types are finite} 


\paragraph{Curry-Howard isomorphism}
We can also apply the implication introduction and implication elimination rules:
\[
    \begin{prooftree}
        \hypo{\Gamma, \sigma \vdash \tau}
        \infer1[$\rightarrow$-I]{\Gamma \vdash \sigma \rightarrow \tau}
    \end{prooftree}
    \qquad
    \begin{prooftree}
        \hypo{\Gamma \vdash \sigma \rightarrow \tau}
        \hypo{\Gamma \vdash \sigma}
        \infer2[$\rightarrow$-E]{\Gamma \vdash \tau}
    \end{prooftree}
\]
